\section{Introducción}
\label{sec:intro}

\paragraph{Contexto industrial}
La lámina de acero plano que llega a un cliente automotriz, de electrodomésticos o de construcción pasa por una
cadena de procesos en la que cada etapa hereda la geometría y el estado mecánico que deja la anterior. En la planta
de Pesquería de Ternium, el primer eslabón de esa cadena para el producto laminado es el \emph{Molino Caliente~4}
(MC4): un planchón recalentado se desbasta y pasa por un tren acabador de siete castillos hasta convertirse en una
banda de pocos milímetros (espesor mediano de \EspesorMediano{}~mm en los rollos de este estudio), que se enfría en la
mesa de salida y se enrolla. La banda no sale perfectamente plana ni de espesor uniforme a lo ancho: suele ser más
gruesa al centro que en las orillas (\emph{corona}) y puede ser más gruesa de un lado que del otro (\emph{cuña}). El
espesor y la corona son los indicadores de desempeño más críticos de la laminación en caliente~\cite{peng2014}.

Varios días después (\DesfaseMedianoH{}~h de mediana en nuestros datos), una parte de esos rollos pasa por el
\emph{Skinpass} de Pesquería, una laminación de temple: una pasada en frío con una reducción muy pequeña, del orden
de 1\,\%~\cite{kijima2008}, y con tensión aplicada. Su propósito es suprimir el punto de cedencia discontinuo que
produce las bandas de Lüders en el conformado posterior, dar el acabado superficial y corregir la forma de la
lámina~\cite{lenard2014}. Para lo último, la palanca principal es el \emph{bending}: la fuerza que dobla los rodillos de
trabajo y cambia el perfil del claro entre ellos y, con ello, cuánto se alarga cada fibra de la lámina a lo ancho. Si
las fibras se alargan de manera desigual quedan esfuerzos residuales que se manifiestan como ondas de orilla o
pandeo central. Por eso la laminación de temple puede corregir una planeza no uniforme que trae el material, pero
también puede revelarla o crearla, y la planeza no uniforme de entrada es una de sus dificultades
reconocidas~\cite{lenard2014}. En la línea estudiada, cada rollo recibe un valor de bending de receta (\emph{preset})
y un lazo de control lo ajusta durante el proceso.

\paragraph{La decisión de Ternium}
Según el Socio Formador, hoy el Skinpass trabaja con una configuración por tipo de material (caliente, decapado o frío) aplicada a todos los
productos de esa familia. La receta única simplifica la operación, pero supone que productos distintos responden de
forma parecida a la misma configuración. Ternium preguntó qué combinación de parámetros del Skinpass se asocia con la
mejor calidad para cada producto y condición de entrada, para decidir si conviene diferenciar recetas por producto.
La decisión es real y tiene un costo: cada receta adicional se diseña, se valida y se mantiene. Contestarla con
historia de planta exige más que mostrar que los productos llegan distinto: exige saber qué pasaría con un producto si
se le cambiara la receta, una predicción contrafactual y no una descripción~\cite{hernan2019}, y mostrar que productos
distintos responden distinto a la misma receta (una interacción), no solo que parten de niveles distintos.

\paragraph{El giro como hallazgo}
Nuestra primera entrega encontró que esa pregunta, tal como está planteada, no puede contestarse con los registros
que existen, y documentar por qué fue el primer resultado del proyecto. Los 24 archivos iniciales describían el MC4,
no el Skinpass (sus cinco columnas del Skinpass venían vacías). Al obtener los resúmenes por rollo del Skinpass de
2025 y ligarlos con el MC4 aparecieron los otros dos obstáculos: el Skinpass no registra planeza de salida, de modo
que no hay un resultado de calidad de forma contra el cual comparar configuraciones, y el bending real no es una
palanca libre, sino la respuesta de un lazo cerrado a la forma de la lámina. Muchos proyectos industriales de
analítica fallan precisamente por no reconocer esa brecha entre la pregunta y los datos y por optimizar un modelo
predictivo como si sus coeficientes fueran efectos~\cite{box1966,shmueli2010}.

Lo que sí registran los datos es cuánto tuvo que corregir el lazo sobre la receta: la diferencia entre el bending
real y el preset, que llamamos $\Y$. Si la receta anticipara lo que el rollo necesita, $\Y$ sería cero. La pregunta
respondible que sigue importando para la decisión es, entonces:
\emph{en los rollos que el MC4 envió al Skinpass de Pesquería en 2025, ¿qué tanto depende la corrección de bending
de lo que se sabe antes de procesarlos ---su producto, su dimensión, el programa del Skinpass y la corona relativa
con que salen del MC4--- y cuál de esas fuentes pesa más?} Aquí, el \emph{programa} es la tanda de rollos que el
Skinpass procesa de corrido (con mediana de \MedianaRollosPrograma{} rollos), y la \emph{corona relativa} es la corona
dividida entre el espesor, en puntos porcentuales (pp). Si la corrección depende sobre todo del producto, una receta
por producto podría absorberla; si depende de otra cosa, diferenciar por producto no la absorbería.

\paragraph{Contribuciones}
Cada una se verifica en la sección~\ref{sec:resultados}:
\begin{enumerate}
\item Una unión verificada, rollo por rollo, entre el MC4 y el Skinpass de un año de producción
      (\NRollosAnalitica{} rollos en \NProgramas{} programas), con la población documentada paso a paso y las
      exclusiones comparadas contra los incluidos.
\item Una descomposición de la corrección de bending contra el azar: el programa del Skinpass explica
      \ExcesoPrograma{}\,\% de su variación por encima de lo que daría el azar, frente a \ExcesoProducto{}\,\% del producto
      y \ExcesoDimension{}\,\% de la dimensión; agrupar por día da casi lo mismo que por programa.
\item Una asociación ajustada entre la corona relativa del MC4 y la corrección, con el signo físico esperado pero
      pequeña (\PendCorona{}~kN por punto porcentual) y frágil ante confusión no medida, sometida a un control negativo,
      a \NVariantes{} alternativas de depuración y agrupación y a un análisis de sensibilidad.
\item Una evaluación predictiva honesta: con validación por programa y temporal, la información previa al Skinpass no
      anticipa la corrección claramente mejor que una receta por producto y dimensión.
\item Una propuesta de experimento evolutivo por programa, con su volumen estimado, y la lista de datos del Skinpass
      que harían respondible la pregunta original.
\end{enumerate}
Mostramos que la corrección que el lazo de bending aplica sobre la receta depende más del programa del Skinpass, y del
periodo en que se procesa el rollo, que del producto, y que la condición de salida del MC4 apenas la anticipa.
